% GENERATED FILE. Edit canonical lessons, then run npm run generate:books.
% canonical-source-hash: fnv1a64:a0d9fd79cf06fd8d
% canonical-lessons: RU-C528-pristavka, RU-C528-videoigra, RU-C528-bloger, RU-C528-podpischik, RU-C528-geymer

\chapter{Consoles and Bloggers}
\label{ch:ru-consoles-and-bloggers}
\hlchaptermodality{528}

\begin{chapteropening}
\textbf{By the end of this chapter:} \emph{I can talk about video games and people online.}

Complete the last lesson of chapter 528: I can talk about video games and people online.
\end{chapteropening}

\section[\texorpdfstring{\ru{приставка} (pristávka)}{приставка (pristávka)}]{\ru{приставка} (pristávka) — a games console; a prefix}

\label{lesson:RU-C528-pristavka}



\begin{quote}
Before the new one: say the Russian for to switch, then the Russian for a screen adaptation.
\end{quote}

\subsection*{You'll want to know: \ru{приставка}}
\textbf{\ru{приставка}} (\emph{pristávka}) — \textquotedblleft{}a games console; a prefix\textquotedblright{}.

\begin{quote}
\ru{На} \ru{день} \ru{рождения} \ru{сын} \ru{получил} \ru{игровую} \ru{приставку}.
\end{quote}

\begin{quote}
\emph{Na den rozhdéniya syn poluchíl igrovúyu pristávku.}
\end{quote}

\begin{quote}
For his birthday my son got a games console.
\end{quote}

Feminine. \emph{\ru{Игровая} \ru{приставка}} — a games console. In grammar, a prefix.

\begin{cousinweb}
\emph{\ru{приставка}} is built on \emph{\ru{приставить}} 'to put next to, to attach', from \emph{\ru{ставить}} 'to stand, to put'; EtymDB lists the word without a source.
\end{cousinweb}

\begin{grammarlens}[title={What you've built}]
One more word for telling, explaining and describing at B1.
\end{grammarlens}

\subsection*{Your turn}
\begin{itemize}
\raggedright
  \item \emph{Say it:} \emph{pristávka}
  \item \emph{Say it:} \emph{pristávka}, once more
  \item \emph{Say it:} a sentence of your own with \emph{pristávka}
  \item \emph{Recall:} say the Russian for to switch, then the Russian for a screen adaptation, then say \emph{pristávka} again
\end{itemize}

\begin{tcolorbox}[breakable,colback=teal!4,colframe=teal!35!black,title={Before you move on}]
What does \ru{приставка} mean? (\textquotedblleft{}A games console; a prefix\textquotedblright{}.) Say it once more.
\end{tcolorbox}
\section[\texorpdfstring{\ru{видеоигра} (videoigrá)}{видеоигра (videoigrá)}]{\ru{видеоигра} (videoigrá) — a video game}

\label{lesson:RU-C528-videoigra}



\begin{quote}
Before the new one: say the Russian for a screen adaptation, then the Russian for a games console.
\end{quote}

\subsection*{You'll want to know: \ru{видеоигра}}
\textbf{\ru{видеоигра}} (\emph{videoigrá}) — \textquotedblleft{}a video game\textquotedblright{}.

\begin{quote}
\ru{Он} \ru{играет} \ru{в} \ru{видеоигры} \ru{по} \ru{выходным}.
\end{quote}

\begin{quote}
\emph{On igráyet v videoígry po vykhodným.}
\end{quote}

\begin{quote}
He plays video games at the weekend.
\end{quote}

Feminine; plural \ru{видеоигры}.

\begin{cousinweb}
\emph{\ru{видеоигра}} is a compound of \emph{\ru{видео}} 'video' and \emph{\ru{игра}} 'a game'; EtymDB has no entry.
\end{cousinweb}

\begin{grammarlens}[title={What you've built}]
One more word for telling, explaining and describing at B1.
\end{grammarlens}

\subsection*{Your turn}
\begin{itemize}
\raggedright
  \item \emph{Say it:} \emph{videoigrá}
  \item \emph{Say it:} \emph{videoigrá}, once more
  \item \emph{Say it:} a sentence of your own with \emph{videoigrá}
  \item \emph{Recall:} say the Russian for a screen adaptation, then the Russian for a games console, then say \emph{videoigrá} again
\end{itemize}

\begin{tcolorbox}[breakable,colback=teal!4,colframe=teal!35!black,title={Before you move on}]
What does \ru{видеоигра} mean? (\textquotedblleft{}A video game\textquotedblright{}.) Say it once more.
\end{tcolorbox}
\section[\texorpdfstring{\ru{блогер} (blóger)}{блогер (blóger)}]{\ru{блогер} (blóger) — a blogger}

\label{lesson:RU-C528-bloger}



\begin{quote}
Before the new one: say the Russian for a games console, then the Russian for a video game.
\end{quote}

\subsection*{You'll want to know: \ru{блогер}}
\textbf{\ru{блогер}} (\emph{blóger}) — \textquotedblleft{}a blogger\textquotedblright{}.

\begin{quote}
\ru{Этот} \ru{блогер} \ru{рассказывает} \ru{о} \ru{путешествиях}.
\end{quote}

\begin{quote}
\emph{Étot blóger rasskázyvayet o puteshéstviyakh.}
\end{quote}

\begin{quote}
This blogger talks about travel.
\end{quote}

Masculine; also of a woman.

\begin{cousinweb}
\emph{\ru{блогер}} is usually traced to English \emph{blogger}; EtymDB has no entry.
\end{cousinweb}

\begin{grammarlens}[title={What you've built}]
One more word for telling, explaining and describing at B1.
\end{grammarlens}

\subsection*{Your turn}
\begin{itemize}
\raggedright
  \item \emph{Say it:} \emph{blóger}
  \item \emph{Say it:} \emph{blóger}, once more
  \item \emph{Say it:} a sentence of your own with \emph{blóger}
  \item \emph{Recall:} say the Russian for a games console, then the Russian for a video game, then say \emph{blóger} again
\end{itemize}

\begin{tcolorbox}[breakable,colback=teal!4,colframe=teal!35!black,title={Before you move on}]
What does \ru{блогер} mean? (\textquotedblleft{}A blogger\textquotedblright{}.) Say it once more.
\end{tcolorbox}
\section[\texorpdfstring{\ru{подписчик} (podpíschik)}{подписчик (podpíschik)}]{\ru{подписчик} (podpíschik) — a subscriber, a follower}

\label{lesson:RU-C528-podpischik}



\begin{quote}
Before the new one: say the Russian for a video game, then the Russian for a blogger.
\end{quote}

\subsection*{You'll want to know: \ru{подписчик}}
\textbf{\ru{подписчик}} (\emph{podpíschik}) — \textquotedblleft{}a subscriber, a follower\textquotedblright{}.

\begin{quote}
\ru{У} \ru{неё} \ru{уже} \ru{миллион} \ru{подписчиков}.
\end{quote}

\begin{quote}
\emph{U neyó uzhé millión podpíschikov.}
\end{quote}

\begin{quote}
She already has a million followers.
\end{quote}

Masculine; feminine \ru{подписчица}.

\begin{cousinweb}
\emph{\ru{подписчик}} is built on \emph{\ru{подписаться}} 'to subscribe', from \emph{\ru{под}-} and \emph{\ru{писать}} 'to write'; EtymDB has no entry.
\end{cousinweb}

\begin{grammarlens}[title={What you've built}]
One more word for telling, explaining and describing at B1.
\end{grammarlens}

\subsection*{Your turn}
\begin{itemize}
\raggedright
  \item \emph{Say it:} \emph{podpíschik}
  \item \emph{Say it:} \emph{podpíschik}, once more
  \item \emph{Say it:} a sentence of your own with \emph{podpíschik}
  \item \emph{Recall:} say the Russian for a video game, then the Russian for a blogger, then say \emph{podpíschik} again
\end{itemize}

\begin{tcolorbox}[breakable,colback=teal!4,colframe=teal!35!black,title={Before you move on}]
What does \ru{подписчик} mean? (\textquotedblleft{}A subscriber, a follower\textquotedblright{}.) Say it once more.
\end{tcolorbox}
\section[\texorpdfstring{\ru{геймер} (géymer)}{геймер (géymer)}]{\ru{геймер} (géymer) — a gamer}

\label{lesson:RU-C528-geymer}



\begin{quote}
Before the new one: say the Russian for a blogger, then the Russian for a subscriber.
\end{quote}

\subsection*{You'll want to know: \ru{геймер}}
\textbf{\ru{геймер}} (\emph{géymer}) — \textquotedblleft{}a gamer\textquotedblright{}.

\begin{quote}
\ru{Мой} \ru{брат} — \ru{настоящий} \ru{геймер}.
\end{quote}

\begin{quote}
\emph{Moy brat — nastoyáshchiy géymer.}
\end{quote}

\begin{quote}
My brother is a real gamer.
\end{quote}

Masculine, colloquial.

\begin{cousinweb}
\emph{\ru{геймер}} is usually traced to English \emph{gamer}; EtymDB has no entry.
\end{cousinweb}

\begin{grammarlens}[title={What you've built}]
One more word for telling, explaining and describing at B1.
\end{grammarlens}

\subsection*{Your turn}
\begin{itemize}
\raggedright
  \item \emph{Say it:} \emph{géymer}
  \item \emph{Say it:} \emph{géymer}, once more
  \item \emph{Say it:} a sentence of your own with \emph{géymer}
  \item \emph{Recall:} say the Russian for a blogger, then the Russian for a subscriber, then say \emph{géymer} again
\end{itemize}

\begin{tcolorbox}[breakable,colback=teal!4,colframe=teal!35!black,title={Before you move on}]
What does \ru{геймер} mean? (\textquotedblleft{}A gamer\textquotedblright{}.) Say it once more.
\end{tcolorbox}
